
\section{Work Plan and Schedule}

In this section we describe which activities we plan on achieving during the time of execution of the scholarship. We also display the timetable for when we expect them to be done and the given priorities for each time period.

\subsection{Planned Activities}

\begin{enumerate}
	\item Bibliographical Research: Praxis in Mathematics, the bibliographic research is expected and needed as a way to get used to important concepts related to the field of research;
	\item Guidance Meetings: As usual for a Ph.D., guidance meetings with the supervisor are both a way of developing the necessary tools of the field and validating the work done so far in the research;
	\item Program's Courses: Expected from any student enrolled in the Mathematics program, there is a quantity of needed course credits one must fulfill;
	\item Qualification Exam: A demand from the program, the qualification exam marks the passage to a Ph.D. candidate;
	\item FAPESP Technical Reports: The technical reports are expected to be completed every year in accordance to the FAPESP guide;
	\item Research Problem Preliminary Study: The first step in the research, along with the bibliographic research, this step marks the beginning of taking the research problem face on;
	\item Research Problem Solution Possibilities: Scavenging the bibliography pops up many possibilities of solutions to the problem, here we endorse some of the possibilities;
	\item Implementation of Possible Solutions: Of course, this step is self-explanatory: One must apply the solutions one has researched;
	\item Results Formulation and Publications: This is one of the final steps in the research process. Formulating the results and publishing are necessary steps for the completion of the Ph.D.;
	\item Thesis Development: The final product of the Ph.D. is, invariably, the thesis. Therefore, time must be reserved for its elaboration;
	\item Theses Defense: The final step in the program, one must defend the thesis one wrote.
\end{enumerate}


\subsection{Execution Schedule}

This P.h.D. research proposal is schedule to start in August 2026 and run for the usual time for the Doctoral Program, that is, 48 months. The project will be run with weekly meetings with the supervisor, and sporadic seminar expositions to the other members of the research team and correlated groups such as the analysis and physical mathematics groups. We can divide the project executions in the following schedule based on the previous steps.
\begin{table}[htbp] \centering
	\label{tab:1}
	\begin{tabular}{l|c|c|c|c|c|c|c|c|}
		\hline
		\multicolumn{ 1}{c|}{\textbf{Activity}} & \multicolumn{8}{c|}{\textbf{Semester}} \\ 
		\cline{ 2- 9}
		\multicolumn{ 1}{l|}{} & \textbf{1} & \textbf{2} & \textbf{3} & \textbf{4} & \textbf{5} & \textbf{6} & \textbf{7} & \textbf{8} \\ \hline
		1 - Bibliographical Research & * & * & \textcolor{red}{*} & \textcolor{red}{*} & * &* & * & *  \\ \hline
		2 - Guidance Meetings & * & * & \textcolor{red}{*} & \textcolor{red}{*} & * & * & * & * \\ \hline
		3 - Program's Courses & \textcolor{red}{*} & \textcolor{red}{*} & * & * &  &   &  & \\ \hline
		4 - Qualification Exam &  \textcolor{red}{*}  & \textcolor{red}{*} & & & & & & \\ \hline
	    5 - FAPESP Technical Reports & & \textcolor{red}{*} & & \textcolor{red}{*} & & \textcolor{red}{*} & & \textcolor{red}{*} \\ \hline
		6 - Research Problem Preliminary Study &  & * &  \textcolor{red}{*}  & & & & & \\ \hline
		7 - Research Problem Solution Possibilities &  & * & \textcolor{red}{*} & \textcolor{red}{*} & & & &  \\ \hline
		8 - Implementation of Possible Solutions & & & & * &* & \textcolor{red}{*} & * & * \\ \hline
		9 - Results Formulation and Publications & & & & & & * & \textcolor{red}{*} & * \\ \hline
		10 - Thesis Development & & & & & & * & \textcolor{red}{*} & * \\ \hline
		11 - Theses Defense & & & & & & & & \textcolor{red}{*} \\ \hline
	\end{tabular}
		\caption{Activities planned, the goals are as described in the Project. Red color indicates the focus of the quarter that may be divided between activities for some crucial quarters.}
\end{table}
